\begin{question}
\noindent A chemistry laboratory uses two similar conical flasks, a small one and a large one, which are mathematically similar in shape.

\noindent The diagram shows the two flasks. The small flask has a height of $12$ cm and a neck width of $3$ cm. The large flask has a height of $30$ cm.

\begin{center}
\begin{tikzpicture}[scale=0.5]
    % Small flask
    \draw[thick] (0,12) -- (-1.5,12) ;
    \draw[thick] (1.5,12) -- (0,12);
    \draw[thick] (-1.5,12) -- (-1.5,10);
    \draw[thick] (1.5,12) -- (1.5,10);
    \draw[thick] (-1.5,10) .. controls (-4,6) and (-4,3) .. (-3,0);
    \draw[thick] (1.5,10) .. controls (4,6) and (4,3) .. (3,0);
    \draw[thick] (-3,0) -- (3,0);
    \draw[<->] (-2.3,10) -- (-2.3,12);
    \node at (-3.6,11) {\footnotesize $3$cm};
    \draw[<->] (5,0) -- (5,12);
    \node at (6.2,6) {\footnotesize $12$cm};

    % Large flask (scaled visually larger, drawn separately to the right)
    \begin{scope}[xshift=15cm, scale=1.6]
    \draw[thick] (0,12) -- (-1.5,12) ;
    \draw[thick] (1.5,12) -- (0,12);
    \draw[thick] (-1.5,12) -- (-1.5,10);
    \draw[thick] (1.5,12) -- (1.5,10);
    \draw[thick] (-1.5,10) .. controls (-4,6) and (-4,3) .. (-3,0);
    \draw[thick] (1.5,10) .. controls (4,6) and (4,3) .. (3,0);
    \draw[thick] (-3,0) -- (3,0);
    \draw[<->] (-2.3,10) -- (-2.3,12);
    \node at (-4.0,11) {\footnotesize $y$cm};
    \draw[<->] (5,0) -- (5,12);
    \node at (6.3,6) {\footnotesize $30$cm};
    \end{scope}
\end{tikzpicture}
\end{center}

\noindent The neck width of the small flask is $y$ cm on the large flask (the flasks are similar, so all corresponding lengths are in the same ratio).

\noindent Work out the value of $y$, the neck width of the large flask.
\vspace{4cm}
\begin{flushright}
\answerequnit{y}{$\text{cm}$}{3}
\end{flushright}
\end{question}